\ExerciseSection

* 1) 设 \(\mathbf{I}\) 是 \(\mathbf{R}^n\) 中的闭立方体，它是 \(n\) 个都等于 \([-\frac{1}{2}\pi, \frac{1}{2}\pi]\) 的区间之积。对每个点 \(x = (x_i)\) （属于 \(\mathbf{I}\) ），我们使它对应于满足下列条件的点 \(y = (y_j)\) （ \(\mathbf{R}^{n+1}\) 中的）：

\[
y _ {1} = \sin x _ {1}
\]

\[
{y _ {2}} {= \cos x _ {1} \sin x _ {2}}
\]

\[
y _ {p} = \cos x _ {1} \cos x _ {2} \dots \cos x _ {p - 1} \sin x _ {p} (2 \leqslant p \leqslant n - 1)
\]

\[
y _ {n} = \cos x _ {1} \cos x _ {2} \dots \cos x _ {n - 1} \sin 2 x _ {n}
\]

\[
y _ {n + 1} = \cos x _ {1} \cos x _ {2} \dots \cos x _ {n - 1} \cos 2 x _ {n}.
\]

证明：I 在这个映射下的象是球面 \(S_{n}\) ，并且这个映射在 I 的内部上的限制是到 \(S_{n}\) 中一个稠密开集上的同胚。 *

\begin{enumerate}
\def\labelenumi{\arabic{enumi})}
\setcounter{enumi}{1}
\tightlist
\item
  定义 \(\mathbf{S}_p \times \mathbf{S}_q\) 到 \(\mathbf{S}_{p + q + 1}\) 的一个子集上的同胚{[}注意 \(\mathbf{S}_{p + q + 1}\) 的方程是
\end{enumerate}

\[
\left. \left(x _ {1} ^ {2} + \dots + x _ {p + 1} ^ {2}\right) + \left(x _ {p + 2} ^ {2} + \dots + x _ {p + q + 2} ^ {2}\right) = \mathrm{I} \right].
\]

\begin{enumerate}
\def\labelenumi{\arabic{enumi})}
\setcounter{enumi}{2}
\tightlist
\item
  \begin{enumerate}
  \def\labelenumii{\alph{enumii})}
  \tightlist
  \item
    证明：若 \(f\) 是 \(\mathbf{S}_1\) 到 \(\mathbf{R}^n\) 中的连续映射，则 \(f\) 可以延拓为 \(\mathbf{B}_2\) 到 \(\mathbf{R}^n\) 中的连续映射{[}把点 \(tx \in \mathbf{B}_2\) （ \(t \geqslant 0\) ，\(x \in S_1\) ）映为 \(tf(x)\) {]}。
  \end{enumerate}
\end{enumerate}

\begin{enumerate}
\def\labelenumi{\alph{enumi})}
\setcounter{enumi}{1}
\tightlist
\item
  由此推出：若 \(f\) 是 \(\mathbf{S}_1\) 到 \(\mathbf{S}_n\) 中的连续映射，满足 \(f(\mathbf{S}_1) \neq \mathbf{S}_n\) ，则 \(f\) 可以延拓为 \(\mathbf{B}_2\) 到 \(\mathbf{S}_n\) 中的连续映射{[}利用一个其顶点不属于 \(f(\mathbf{S}_1)\) 的球极投影{]}。(*)
\end{enumerate}

\begin{enumerate}
\def\labelenumi{\arabic{enumi})}
\setcounter{enumi}{3}
\tightlist
\item
  证明：不存在 \(S_{1}\) 到 R 中的同胚。（注意 \(S_{1}\) 中任一点在 \(S_{1}\) 中的补集是连通的，而 R 的每个连通子集是一个区间。）
\end{enumerate}

由此推出：\(S_{1}\) 到 \(S_{1}\) 的一个子空间上的同胚必然是 \(S_{1}\) 到 \(S_{1}\) 上的同胚（利用第 4 小节的命题 4）。

\begin{enumerate}
\def\labelenumi{\arabic{enumi})}
\setcounter{enumi}{4}
\item
  证明：若 \(n > 1\) ，则球面 \(\mathbf{S}_n\) 不同胚于圆周 \(\mathbf{S}_1\) （参见 § 1 的习题 8）。
\item
  把 \(S_{1}\) 与环面 T（即 R 关于关系 \(x \equiv y \pmod{1}\) 的商，第 4 小节命题 4 的推论 2）等同起来，设 \(\varphi\) 表示 R 到 \(S_{1}\) 上的典范映射。拓扑空间 E 到 \(S_{1}\) 中的连续映射 f 称为非本质的，如果存在 E 到 R 中的连续映射 g，使得 \(f = \varphi \circ g\) 。(**) 不是非本质的映射称为本质映射。
\end{enumerate}

证明 \(S_{1}\) 到 \(S_{1}\) 上的恒等映射是本质映射（利用习题 4）。

\begin{enumerate}
\def\labelenumi{\arabic{enumi})}
\setcounter{enumi}{6}
\item
  证明：存在 \(S_{1}\) 的一致结构的一个围域 U，使得只要 f 是拓扑空间 E 到 \(S_{1}\) 中的非本质映射，则每个连续映射 \(f' : E \to S_{1}\) （它使 \((f(x), f'(x)) \in U\) 对一切 \(x \in E\) 成立）也是非本质的。
\item
  证明：不存在连续映射 \(f\) ：它把 \(\mathbf{B}_2\) 映到 \(\mathbf{S}_1\) 上并延拓 \(\mathbf{S}_1\) 的恒等映射（利用习题 7，证明 \(f\) 限制在以 o 为中心、以 \(r \leqslant i\) 为半径的圆周上时，总是该圆周到 \(S_{1}\) 中的非本质映射，从而当 \(r = i\) 时得出矛盾）。
\item
  设 E 是含多于一个点的拓扑空间。证明：存在 \(S_{1}\) 到 \(F = E \times S_{1}\) 中的连续映射 f，满足 \(f(S_{1}) \neq F\) ，并且它不能延拓为 \(B_{2}\) 到 F 中的连续映射（利用习题 8）。由此推出：若 n \textgreater{} 1 ，则 \(S_{n}\) 、\(R^{n}\) 与 \(B_{n}\) 都不同胚于任何形如 \(E \times S_{1}\) 的空间，其中 E 是任意拓扑空间（参见习题 3）。特别地，若 n \textgreater{} 1 ，\(S_{n}\) 不同胚于 \(S_{1}^{n}\) ，而且对一切 n，\(B_{n}\) 不同胚于 \(S_{1}^{n}\) 。
\end{enumerate}

同样证明：\(R^{2}\) 不同胚于 \(R^{2}\) 中一点的补集，且 \(S_{2}\) 不同胚于 \(B_{2}\) （若不真，\(R^{2}\) 将同胚于 \(S_{1}\) 与区间 \([0, 1]\) 的乘积）。

\begin{enumerate}
\def\labelenumi{\arabic{enumi})}
\setcounter{enumi}{9}
\tightlist
\item
  设 \(H_{n,p,q}\) 表示 \(R^{n}\) 中由方程
\end{enumerate}

\[
x _ {1} ^ {2} + x _ {2} ^ {2} + \dots + x _ {p} ^ {2} - x _ {p + 1} ^ {2} - \dots - x _ {p + q} ^ {2} = \mathrm{I} \quad (p + q \leqslant n).
\]

所定义的“二次超曲面”。证明 \(\mathbf{H}_{n,p,q}\) 同胚于 \(\mathbf{S}_{p - 1}\times \mathbf{R}^{n - p}\) 。

\begin{enumerate}
\def\labelenumi{\arabic{enumi})}
\setcounter{enumi}{10}
\tightlist
\item
  设 \(C_{n,p}\) 是 \(R^{n}\) 中由方程
\end{enumerate}

\[
x _ {1} ^ {2} + x _ {2} ^ {2} + \dots + x _ {p} ^ {2} - x _ {p + 1} ^ {2} - \dots - x _ {n} ^ {2} = 0 \quad (\mathrm{I} \leqslant p \leqslant n - \mathrm{I}).
\]

所定义的“二次锥面”，证明 \(\{0\}\) 在 \(\mathbf{C}_{n,p}\) 中的补集同胚于 \(\mathbf{S}_{p-1} \times \mathbf{S}_{n-p-1} \times \mathbf{R}\) 。

* 12) \(\mathbf{R}^n\) 中含原点 o 的子集 E 称为（关于 o 的）星形集，如果对一切 \(x \in \mathbf{E}\) 与一切 \(t \in [0, 1]\) ，都有 \(tx \in \mathbf{E}\) 。E 与以 o 为端点的闭射线的交或者是整条射线，或者是一条以 o 为其一个端点的线段。E 的壳是指由这些线段的非零端点组成的集 K，再加上 o——如果存在一条射线，它与 E 的交仅由 o 组成。在下文中我们总假设 \(o \notin K\) 。

\begin{enumerate}
\def\labelenumi{\alph{enumi})}
\item
  证明 E 的壳含于 E 的边界。给出一个这两个集不相同的例子。
\item
  证明：若 E 的壳 K 是紧的，则存在 \(R^{n}\) 到其自身的同胚，它把 K 变换成 \(S_{n-1}\) ，把 \(\overline{E}\) 变换成 \(B_{n}\) ，并且把 E 的内部变换成 \(B_{n}\) 的内部（把每一点 \(x \in K\) 映为它在 \(S_{n-1}\) 上的中心投影，然后把这个映射延拓到整个 \(R^{n}\) 上）。由此推出：E 的边界与其壳重合，并且 E 的内部是由点 tx 组成的集，其中 \(x \in K\) 且 \(t \in [0, I[\) 。
\item
  若 E 无界且其边界与其壳重合，则 E 的内部同胚于 \(R^{n}\) ，并且它的壳 K 同胚于 \(S_{n-1}\) 的一个开子集{[}证明 E 在同胚 \(x \to x/(1 + ||x||)\) 下的象满足 b) 的条件{]}。
\item
  给出一个无界星形集 E 的例子，它的壳是闭的，但不与 E 的边界重合。
\end{enumerate}

\begin{enumerate}
\def\labelenumi{\arabic{enumi})}
\setcounter{enumi}{12}
\tightlist
\item
  证明：在空间 \(S_{n}\) 中，外部非空的非空开集的边界具有连续统的势（利用第 4 小节的命题 4）。
\end{enumerate}

